
\documentclass[3pt]{article}

\begin{document}
	
	DECISION TREES — ADVANCED EXAM WITH ANSWERS
	
\vspace{0.3em}

	==================================================
	
	PARAGRAPH 1 — DECISION TREE BASICS
	
	\begin{enumerate}
		\item A decision tree is mainly used for:
		\begin{enumerate}
			\item Clustering
			\item Classification
			\item Dimensionality reduction
			\item Optimization
		\end{enumerate}
		
		\item In a decision tree, internal nodes represent:
		\begin{enumerate}
			\item Class labels
			\item Attribute tests
			\item Training instances
			\item Probabilities
		\end{enumerate}
		
		\item Leaves of a decision tree represent:
		\begin{enumerate}
			\item Attributes
			\item Thresholds
			\item Class labels
			\item Features
		\end{enumerate}
		
		\item Branches leaving a node correspond to:
		\begin{enumerate}
			\item Classes
			\item Attribute values
			\item Weights
			\item Errors
		\end{enumerate}
		
		\item Decision trees are learned using:
		\begin{enumerate}
			\item Unsupervised learning
			\item Supervised learning
			\item Reinforcement learning
			\item Self-supervised learning
		\end{enumerate}
		
		\item Traversing a tree means:
		\begin{enumerate}
			\item Updating weights
			\item Evaluating attributes sequentially
			\item Computing probabilities
			\item Optimizing loss
		\end{enumerate}
		
		\item A decision tree classifier is:
		\begin{enumerate}
			\item Probabilistic
			\item Deterministic
			\item Stochastic
			\item Randomized
		\end{enumerate}
		
		\item Each internal node usually tests:
		\begin{enumerate}
			\item All attributes
			\item One attribute
			\item All classes
			\item One instance
		\end{enumerate}
		
		\item Decision trees can be interpreted as:
		\begin{enumerate}
			\item Black-box models
			\item Rule-based models
			\item Neural networks
			\item Kernel machines
		\end{enumerate}
		
		\item A main advantage of decision trees is:
		\begin{enumerate}
			\item High interpretability
			\item Low variance
			\item Convex optimization
			\item Probabilistic outputs
		\end{enumerate}
	\end{enumerate}
	
	Correct answers: 1B, 2B, 3C, 4B, 5B, 6B, 7B, 8B, 9B, 10A
	
\vspace{0.3em}

	==================================================
	
	PARAGRAPH 2 — LOGICAL INTERPRETATION (RULES)
	
	\begin{enumerate}
		\item A root-to-leaf path represents:
		\begin{enumerate}
			\item A disjunction
			\item A conjunction
			\item A negation
			\item A probability
		\end{enumerate}
		
		\item A complete decision tree represents:
		\begin{enumerate}
			\item A conjunction of disjunctions
			\item A disjunction of conjunctions
			\item A single rule
			\item A linear classifier
		\end{enumerate}
		
		\item Each rule extracted from a tree corresponds to:
		\begin{enumerate}
			\item One attribute
			\item One branch
			\item One path
			\item One node
		\end{enumerate}
		
		\item Decision tree rules are usually:
		\begin{enumerate}
			\item Probabilistic
			\item Fuzzy
			\item Deterministic
			\item Approximate
		\end{enumerate}
		
		\item If an instance satisfies multiple rules:
		\begin{enumerate}
			\item A conflict occurs
			\item One leaf is reached
			\item Probabilities are averaged
			\item Classification fails
		\end{enumerate}
		
		\item Rule extraction improves:
		\begin{enumerate}
			\item Accuracy
			\item Interpretability
			\item Training speed
			\item Generalization
		\end{enumerate}
		
		\item Logical form of trees is typically:
		\begin{enumerate}
			\item CNF
			\item DNF
			\item Linear
			\item Polynomial
		\end{enumerate}
		
		\item Each rule ends with:
		\begin{enumerate}
			\item Attribute test
			\item Class assignment
			\item Probability estimate
			\item Threshold
		\end{enumerate}
		
		\item Removing a rule corresponds to:
		\begin{enumerate}
			\item Removing a leaf
			\item Removing a root
			\item Removing an attribute
			\item Removing all branches
		\end{enumerate}
		
		\item Rule-based form is useful for:
		\begin{enumerate}
			\item Model explanation
			\item Kernel design
			\item Gradient computation
			\item Dimensionality reduction
		\end{enumerate}
	\end{enumerate}
	
	Correct answers: 1B, 2B, 3C, 4C, 5B, 6B, 7B, 8B, 9A, 10A
	
\vspace{0.3em}

	==================================================
	
	PARAGRAPH 3 — SPACE PARTITIONING
	
	\begin{enumerate}
		\item Decision trees partition the feature space into:
		\begin{enumerate}
			\item Spheres
			\item Hyperplanes
			\item Hyperrectangles
			\item Clusters
		\end{enumerate}
		
		\item Each leaf corresponds to:
		\begin{enumerate}
			\item One point
			\item One region
			\item One attribute
			\item One probability
		\end{enumerate}
		
		\item Decision boundaries of trees are:
		\begin{enumerate}
			\item Linear
			\item Oblique
			\item Axis-aligned
			\item Circular
		\end{enumerate}
		
		\item Axis-aligned splits test:
		\begin{enumerate}
			\item Linear combinations
			\item Single attributes
			\item All attributes
			\item Class labels
		\end{enumerate}
		
		\item Increasing tree depth causes:
		\begin{enumerate}
			\item Fewer regions
			\item Larger regions
			\item More regions
			\item No change
		\end{enumerate}
		
		\item Each split reduces:
		\begin{enumerate}
			\item Feature dimension
			\item Instance count per node
			\item Number of attributes globally
			\item Label space
		\end{enumerate}
		
		\item Regions created by trees are:
		\begin{enumerate}
			\item Overlapping
			\item Disjoint
			\item Probabilistic
			\item Fuzzy
		\end{enumerate}
		
		\item Decision trees can model:
		\begin{enumerate}
			\item Linear boundaries only
			\item Arbitrary shapes via partitions
			\item Circular regions directly
			\item Kernelized spaces
		\end{enumerate}
		
		\item Space partitioning is:
		\begin{enumerate}
			\item Global
			\item Local and recursive
			\item Random
			\item Gradient-based
		\end{enumerate}
		
		\item A shallow tree yields:
		\begin{enumerate}
			\item Complex regions
			\item Simple regions
			\item High variance
			\item Overfitting
		\end{enumerate}
	\end{enumerate}
	
	Correct answers: 1C, 2B, 3C, 4B, 5C, 6B, 7B, 8B, 9B, 10B
	
\vspace{0.3em}

	==================================================
	
	PARAGRAPH 4 — TREE GENERATION ALGORITHMS
	
	\begin{enumerate}
		\item ID3 is an algorithm for:
		\begin{enumerate}
			\item Regression trees
			\item Classification trees
			\item Clustering
			\item Reinforcement learning
		\end{enumerate}
		
		\item ID3 builds trees:
		\begin{enumerate}
			\item Bottom-up
			\item Top-down
			\item Randomly
			\item Backtracking
		\end{enumerate}
		
		\item At each node, ID3 selects:
		\begin{enumerate}
			\item Random attribute
			\item Attribute with highest information gain
			\item Attribute with lowest variance
			\item Attribute with most values
		\end{enumerate}
		
		\item ID3 stops splitting when:
		\begin{enumerate}
			\item Accuracy is maximum
			\item All attributes are used
			\item Dataset is empty or pure
			\item Depth is fixed
		\end{enumerate}
		
		\item After splitting, the chosen attribute is:
		\begin{enumerate}
			\item Reused
			\item Ignored temporarily
			\item Removed from branch
			\item Randomized
		\end{enumerate}
		
		\item ID3 handles continuous attributes:
		\begin{enumerate}
			\item Directly
			\item With thresholds
			\item Not natively
			\item Probabilistically
		\end{enumerate}
		
		\item ID3 was proposed by:
		\begin{enumerate}
			\item Breiman
			\item Quinlan
			\item Vapnik
			\item Hastie
		\end{enumerate}
		
		\item ID3 evaluates attributes using:
		\begin{enumerate}
			\item Gini index
			\item Entropy
			\item Variance
			\item Error rate
		\end{enumerate}
		
		\item ID3 produces trees that tend to:
		\begin{enumerate}
			\item Underfit
			\item Overfit
			\item Be optimal
			\item Be shallow
		\end{enumerate}
		
		\item ID3 assumes attributes are:
		\begin{enumerate}
			\item Continuous
			\item Binary only
			\item Discrete
			\item Ordered
		\end{enumerate}
	\end{enumerate}
	
	Correct answers: 1B, 2B, 3B, 4C, 5C, 6C, 7B, 8B, 9B, 10C
	
\vspace{0.3em}

	==================================================
	
	PARAGRAPH 5 — ENTROPY AND INFORMATION GAIN
	
	\begin{enumerate}
		\item Entropy measures:
		\begin{enumerate}
			\item Accuracy
			\item Uncertainty
			\item Error rate
			\item Variance
		\end{enumerate}
		
		\item Entropy is minimum when:
		\begin{enumerate}
			\item Classes are balanced
			\item One class dominates
			\item Attributes are many
			\item Samples are few
		\end{enumerate}
		
		\item Entropy is maximum when:
		\begin{enumerate}
			\item One class exists
			\item Classes are equally distributed
			\item No attributes exist
			\item Tree is shallow
		\end{enumerate}
		
		\item Information gain measures:
		\begin{enumerate}
			\item Remaining entropy
			\item Reduction in entropy
			\item Classification accuracy
			\item Tree depth
		\end{enumerate}
		
		\item An attribute with high gain:
		\begin{enumerate}
			\item Increases uncertainty
			\item Reduces uncertainty
			\item Is random
			\item Has many values
		\end{enumerate}
		
		\item Weighted entropy accounts for:
		\begin{enumerate}
			\item Tree depth
			\item Sample proportions
			\item Attribute count
			\item Label count
		\end{enumerate}
		
		\item Gain is computed as:
		\begin{enumerate}
			\item Entropy minus weighted entropy
			\item Weighted entropy minus entropy
			\item Sum of entropies
			\item Average entropy
		\end{enumerate}
		
		\item Attributes with many values may:
		\begin{enumerate}
			\item Be penalized
			\item Be favored
			\item Be ignored
			\item Have zero gain
		\end{enumerate}
		
		\item Entropy uses logarithm base:
		\begin{enumerate}
			\item 10
			\item e
			\item 2
			\item n
		\end{enumerate}
		
		\item Information gain is used to:
		\begin{enumerate}
			\item Prune trees
			\item Select split attributes
			\item Compute probabilities
			\item Normalize data
		\end{enumerate}
	\end{enumerate}
	
	Correct answers: 1B, 2B, 3B, 4B, 5B, 6B, 7A, 8B, 9C, 10B
	
\vspace{0.3em}

	==================================================
	
	PARAGRAPH 6 — OVERFITTING AND PRUNING
	
	\begin{enumerate}
		\item Overfitting occurs when:
		\begin{enumerate}
			\item Tree is too shallow
			\item Tree memorizes noise
			\item Data is balanced
			\item Entropy is low
		\end{enumerate}
		
		\item ID3 is prone to overfitting because:
		\begin{enumerate}
			\item It stops early
			\item It grows until purity
			\item It uses gain
			\item It ignores labels
		\end{enumerate}
		
		\item Pre-pruning means:
		\begin{enumerate}
			\item Cutting after training
			\item Stopping early
			\item Removing attributes
			\item Random splitting
		\end{enumerate}
		
		\item Post-pruning means:
		\begin{enumerate}
			\item Stopping early
			\item Removing branches after growth
			\item Adding branches
			\item Changing attributes
		\end{enumerate}
		
		\item Pruning improves:
		\begin{enumerate}
			\item Training accuracy
			\item Interpretability
			\item Generalization
			\item Tree depth
		\end{enumerate}
		
		\item Noise mainly affects:
		\begin{enumerate}
			\item Root node
			\item Deep leaves
			\item Attribute selection
			\item Class labels only
		\end{enumerate}
		
		\item Very deep trees usually have:
		\begin{enumerate}
			\item High bias
			\item High variance
			\item Low variance
			\item No variance
		\end{enumerate}
		
		\item Pruning trades off:
		\begin{enumerate}
			\item Bias and variance
			\item Speed and memory
			\item Accuracy and entropy
			\item Attributes and labels
		\end{enumerate}
		
		\item Pruned trees are generally:
		\begin{enumerate}
			\item Less interpretable
			\item More interpretable
			\item Random
			\item Unstable
		\end{enumerate}
		
		\item The goal of pruning is to:
		\begin{enumerate}
			\item Maximize depth
			\item Minimize error on training data
			\item Improve performance on unseen data
			\item Eliminate entropy
		\end{enumerate}
	\end{enumerate}
	
	Correct answers: 1B, 2B, 3B, 4B, 5C, 6B, 7B, 8A, 9B, 10C
	
\end{document}
